% Source: https://artofproblemsolving.com/community/c3h2512655_infinity
% Stable identity: aops-topic-2512655-q01
Let $f:\mathbb{N}\to\mathcal{P}(\mathbb{N})$ be any function; each $f(n)$ is a subset of $\mathbb{N}$. Define
\[D=\{n\in\mathbb{N}:\ n\notin f(n)\}.\]
This is a subset of $\mathbb{N}$. We show that $D$ is not a value of $f$. Suppose $D=f(k)$ for some $k$.
\begin{itemize}
\item If $k\in D$, then by the definition of $D$, $k\notin f(k)=D$, a contradiction.
\item If $k\notin D$, then $k\notin f(k)$, so by the definition of $D$, $k\in D$, a contradiction.
\end{itemize}
Either way we get a contradiction, so $D\ne f(k)$ for every $k$, and $f$ is not onto. $\blacksquare$

(This is Cantor's diagonal argument: $D$ differs from the set $f(n)$ in whether it contains $n$. Repeating the construction gives larger and larger infinite sets, so there is no largest infinity.)
